% ============================================================================
%  estilo.tex — identidade visual, tipografia e componentes do documento
%
%  Tipografia: Roboto (texto e titulos) + Roboto Mono (codigo). As duas vem do
%  pacote `roboto`, entao o documento compila com pdfLaTeX, XeLaTeX ou Tectonic
%  sem fontspec explicito.
%
%  Paleta: os tokens `--unooria-*` do @unooria/design-system, os MESMOS valores
%  que `packages/ui/styles/plataforma.css` usa como fallback. Cor nova neste
%  documento so entra se existir como token no pacote.
% ============================================================================

% --- Codificacao e idioma ---------------------------------------------------
\usepackage[T1]{fontenc}
\usepackage[utf8]{inputenc}
\usepackage[shorthands=off,brazil]{babel}   % shorthands OFF: com babel-pt o `"` fica ativo e quebra codigo em linha
\usepackage{csquotes}

% --- Tipografia -------------------------------------------------------------
\usepackage[sfdefault]{roboto}            % Roboto como familia padrao
\usepackage[scaled=0.90]{roboto-mono}     % Roboto Mono para codigo
\usepackage{microtype}
\usepackage{textcomp}

% --- Pagina -----------------------------------------------------------------
\usepackage[
  a4paper,
  top=2.6cm, bottom=2.5cm, left=2.6cm, right=2.4cm,
  headheight=15pt, headsep=18pt, footskip=26pt
]{geometry}
\usepackage{pdflscape}
\usepackage{setspace}
\onehalfspacing

% --- Cores canônicas do Brandbook Unooria ("BLUNOORIA") -----------------------
% Regra semântica: "A cor não decora. Ela organiza."
%   - Azul Profundo (#0E1039): Estrutura, corpo, regras e títulos
%   - Lavanda (#8B85AC): Inteligência, decisões arquiteturais (ADR), IA
%   - Azul Claro (#7FA5E2): Interface, navegação, referências e notas
\usepackage{xcolor}
\definecolor{unooNavy}        {HTML}{0E1039}  % Azul Profundo — Estrutura (C100 M96 Y44 K56)
\definecolor{unooLavender}    {HTML}{8B85AC}  % Lavanda — Inteligência (C52 M48 Y15 K2)
\definecolor{unooBlue}        {HTML}{7FA5E2}  % Azul Claro — Interface (C54 M29 Y0 K0)
\definecolor{unooBlueLight}   {HTML}{B8D0F5}  % Variação clara de interface
\definecolor{unooBlueWash}    {HTML}{F0F4FC}  % Fundo suave de interface
\definecolor{unooLavenderWash}{HTML}{F3F2F8}  % Fundo suave de inteligência
\definecolor{unooLavenderSoft}{HTML}{5D577D}  % Lavanda escurecido p/ alto contraste tipográfico
\definecolor{unooInk}         {HTML}{0E1039}  % neutral-950 / texto primário ancorado na marca
\definecolor{unooInkSoft}     {HTML}{3A3C59}  % neutral-700 / texto secundário
\definecolor{unooMuted}       {HTML}{6D6E82}  % neutral-500 / legendas e metadados
\definecolor{unooGrid}        {HTML}{E3E5EE}  % Linhas de grade e bordas sutis
\definecolor{unooAxis}        {HTML}{CACDDC}  % Réguas e separadores
\definecolor{unooSurface}     {HTML}{FFFFFF}  % Fundo puro
\definecolor{unooPaper}       {HTML}{F7F8FC}  % Fundo alternado / zebra
% Tons funcionais (alertas de sistema e status):
\definecolor{unooRed}         {HTML}{C92A2A}  % danger-600
\definecolor{unooRedWash}     {HTML}{FDF2F2}
\definecolor{unooOrange}      {HTML}{D97706}  % warning-600
\definecolor{unooOrangeWash}  {HTML}{FFFBEB}
\definecolor{unooGreen}       {HTML}{2F855A}  % Status funcional aprovado
\definecolor{unooGreenWash}   {HTML}{F0FDF4}
\definecolor{unooYellow}      {HTML}{B45309}
\definecolor{unooYellowWash}  {HTML}{FEF3C7}
\definecolor{unooPink}        {HTML}{9D4EDD}

% --- Tabelas ----------------------------------------------------------------
\usepackage{booktabs}
\usepackage{tabularx}
\usepackage{longtable}
\usepackage{array}
\usepackage{multirow}
\usepackage{makecell}
\usepackage{colortbl}
\usepackage{ltablex}          % tabularx com quebra de pagina
\keepXColumns
\renewcommand{\arraystretch}{1.22}
\arrayrulecolor{unooAxis}
\setlength{\aboverulesep}{0pt}
\setlength{\belowrulesep}{0pt}
\setlength{\extrarowheight}{2pt}

% Tabelas em \small por decisao de composicao: as celulas deste documento
% misturam prosa e identificadores monoespacados longos (caminhos de arquivo,
% nomes de coluna), que NAO hifenizam. Em corpo 11 elas transbordavam a caixa;
% em \small cabem, e a tabela continua legivel porque a densidade e alta de
% proposito. `\emergencystretch` cobre o residuo.
\usepackage{etoolbox}
\AtBeginEnvironment{tabularx}{\small\setlength{\emergencystretch}{1.2em}}
\AtBeginEnvironment{longtable}{\small\setlength{\emergencystretch}{1.2em}}

\newcolumntype{L}[1]{>{\raggedright\arraybackslash}p{#1}}
\newcolumntype{C}[1]{>{\centering\arraybackslash}p{#1}}
\newcolumntype{R}[1]{>{\raggedleft\arraybackslash}p{#1}}
\newcolumntype{Y}{>{\raggedright\arraybackslash}X}

% Colunas X sem justificacao: as celulas misturam prosa e codigo monoespacado,
% que nao hifeniza — justificar produziria espacos internos muito largos.
\renewcommand{\tabularxcolumn}[1]{>{\raggedright\arraybackslash}p{#1}}

\newcommand{\thd}[1]{\textcolor{white}{\textbf{#1}}}
\newcommand{\headrow}{\rowcolor{unooNavy}}
\newcommand{\zebra}{\rowcolor{unooPaper}}

% --- Listas -----------------------------------------------------------------
\usepackage{enumitem}
\setlist[itemize]{leftmargin=1.35em, itemsep=2pt, topsep=4pt, parsep=0pt,
                  label=\textcolor{unooBlue}{\raisebox{0.18ex}{\rule{0.40em}{0.40em}}}}
\setlist[enumerate]{leftmargin=1.7em, itemsep=2pt, topsep=4pt, parsep=0pt,
                    font=\bfseries\color{unooNavy}}
\setlist[description]{leftmargin=1.4em, style=nextline,
                      font=\bfseries\color{unooNavy}}

% --- Matematica e unidades --------------------------------------------------
\usepackage{amsmath}
\usepackage{amssymb}
\usepackage{siunitx}
\sisetup{output-decimal-marker={,}, group-separator={.}, group-minimum-digits=4,
         per-mode=symbol, detect-weight=true, detect-family=true}
\DeclareSIUnit{\tokens}{tokens}
\DeclareSIUnit{\req}{req}

% --- Figuras e legendas -----------------------------------------------------
\usepackage{graphicx}
\usepackage{float}
\usepackage{adjustbox}
\usepackage[font=small, labelfont={bf,color=unooNavy}, textfont=normalfont,
            labelsep=endash, justification=raggedright, singlelinecheck=false,
            skip=8pt]{caption}
\usepackage{subcaption}
\captionsetup[table]{position=top, skip=6pt}

% Ajusta a figura a largura do texto apenas quando ela excede a caixa.
\newcommand{\fitwidth}[1]{\adjustbox{max width=\linewidth, center}{#1}}
\newcommand{\fitpage}[1]{\adjustbox{max width=\linewidth, max height=0.82\textheight, center}{#1}}

% --- Blocos de destaque (tcolorbox) ----------------------------------------
\usepackage[most]{tcolorbox}
\tcbuselibrary{skins,breakable}

\tcbset{
  unoobase/.style={
    enhanced, breakable, sharp corners,
    boxrule=0pt, leftrule=2.6pt, arc=0pt,
    left=10pt, right=10pt, top=7pt, bottom=7pt,
    fonttitle=\bfseries\small,
    before skip=10pt, after skip=10pt
  }
}

% `colbacktitle` e OBRIGATORIO aqui: sem ele o tcolorbox pinta a faixa do
% titulo com a cor de `colframe`, e o titulo (que usa a MESMA cor em
% `coltitle`) fica invisivel — texto vermelho sobre faixa vermelha. O sintoma
% e uma barra colorida solida no lugar do rotulo do bloco.
\newtcolorbox{nota}[1][]{unoobase, colback=unooBlueWash, colframe=unooBlue,
  colbacktitle=unooBlueWash, coltitle=unooNavy, titlerule=0pt, title={NOTA}, #1}
\newtcolorbox{atencao}[1][]{unoobase, colback=unooOrangeWash, colframe=unooOrange,
  colbacktitle=unooOrangeWash, coltitle=unooOrange, titlerule=0pt,
  title={ATENÇÃO}, #1}
\newtcolorbox{decisao}[1][]{unoobase, colback=unooLavenderWash, colframe=unooLavender,
  colbacktitle=unooLavenderWash, coltitle=unooNavy, titlerule=0pt,
  title={DECISÃO DE PROJETO}, #1}
\newtcolorbox{limitacao}[1][]{unoobase, colback=unooRedWash, colframe=unooRed,
  colbacktitle=unooRedWash, coltitle=unooRed, titlerule=0pt,
  title={LIMITAÇÃO CONHECIDA}, #1}
\newtcolorbox{regra}[1][]{unoobase, colback=unooRedWash, colframe=unooRed,
  colbacktitle=unooRedWash, coltitle=unooRed, titlerule=0pt,
  title={REGRA INVIOLÁVEL}, #1}
\newtcolorbox{destaque}[1][]{unoobase, colback=unooPaper, colframe=unooLavender, #1}
\newtcolorbox{especificacao}[1]{unoobase, colback=unooSurface, colframe=unooNavy,
  colbacktitle=unooSurface, coltitle=unooNavy, titlerule=0pt, title={#1}}

% --- Codigo-fonte -----------------------------------------------------------
\usepackage{listings}

\definecolor{lstBg}     {HTML}{FAFBFB}
\definecolor{lstFrame}  {HTML}{E4E9EB}
\definecolor{lstKeyword}{HTML}{2F4F7F}
\definecolor{lstString} {HTML}{0F6B52}
\definecolor{lstComment}{HTML}{617379}
\definecolor{lstEmph}   {HTML}{A9490F}

\lstdefinestyle{unoo}{
  backgroundcolor=\color{lstBg},
  basicstyle=\ttfamily\scriptsize\color{unooInk},
  keywordstyle=\color{lstKeyword}\bfseries,
  keywordstyle=[2]\color{lstEmph}\bfseries,
  keywordstyle=[3]\color{unooBlue},
  stringstyle=\color{lstString},
  commentstyle=\color{lstComment}\itshape,
  numberstyle=\ttfamily\tiny\color{unooMuted},
  identifierstyle=\color{unooInk},
  numbers=left, numbersep=9pt, stepnumber=1,
  frame=single, framesep=6pt, rulecolor=\color{lstFrame},
  xleftmargin=17pt, xrightmargin=2pt,
  breaklines=true, breakatwhitespace=false, breakindent=0pt,
  postbreak=\mbox{\textcolor{unooBlue}{$\hookrightarrow$}\space},
  showstringspaces=false, tabsize=2,
  columns=fullflexible, keepspaces=true, upquote=true,
  aboveskip=10pt, belowskip=10pt,
  literate=
    {á}{{\'a}}1 {à}{{\`a}}1 {ã}{{\~a}}1 {â}{{\^a}}1
    {é}{{\'e}}1 {ê}{{\^e}}1 {í}{{\'i}}1
    {ó}{{\'o}}1 {õ}{{\~o}}1 {ô}{{\^o}}1
    {ú}{{\'u}}1 {ü}{{\"u}}1 {ç}{{\c c}}1
    {Á}{{\'A}}1 {À}{{\`A}}1 {Ã}{{\~A}}1 {Â}{{\^A}}1
    {É}{{\'E}}1 {Ê}{{\^E}}1 {Í}{{\'I}}1
    {Ó}{{\'O}}1 {Õ}{{\~O}}1 {Ô}{{\^O}}1
    {Ú}{{\'U}}1 {Ç}{{\c C}}1
    {°}{{\textdegree}}1 {º}{{\textordmasculine}}1
    {µ}{{$\mu$}}1 {×}{{$\times$}}1 {≈}{{$\approx$}}1
    {≥}{{$\geq$}}1 {≤}{{$\leq$}}1 {≠}{{$\neq$}}1
    {→}{{$\rightarrow$}}1 {←}{{$\leftarrow$}}1 {⇒}{{$\Rightarrow$}}1
    {—}{{---}}3 {–}{{--}}2 {…}{{\ldots}}3
    {“}{{\textquotedblleft}}1 {”}{{\textquotedblright}}1
    {‘}{{\textquoteleft}}1 {’}{{\textquoteright}}1
    {ª}{{\textordfeminine}}1 {§}{{\S}}1
}
\lstset{style=unoo}

\lstdefinelanguage{tsunoo}[]{Java}{
  morekeywords={const,let,export,default,import,from,as,interface,type,await,
                async,function,return,string,number,boolean,readonly,unknown,
                never,satisfies,enum,declare,namespace,implements},
  morekeywords=[2]{z,Prisma,useEffect,postMessage},
  sensitive=true,
  morecomment=[l]{//},
  morecomment=[s]{/*}{*/},
  morestring=[b]",
  morestring=[b]',
  morestring=[b]`
}
\lstdefinelanguage{sqlunoo}[]{SQL}{
  morekeywords={POLICY,ENABLE,FORCE,ROW,LEVEL,SECURITY,USING,PARTITION,RANGE,
                CONFLICT,DO,NOTHING,JSONB,TIMESTAMPTZ,current_setting,
                set_config,ATTACH,DEFAULT,GENERATED,ALWAYS},
  sensitive=false,
  morecomment=[l]{--},
  morestring=[b]'
}
\lstdefinelanguage{yamlunoo}{
  morekeywords={true,false,null,services,image,build,command,environment,
                volumes,ports,depends_on,healthcheck,condition,runs-on,steps,
                uses,needs,jobs,name,on,with,env,if,run,concurrency,
                permissions,workflow_call,workflow_dispatch},
  sensitive=true,
  morecomment=[l]{\#},
  morestring=[b]",
  morestring=[b]'
}
\lstdefinelanguage{jsonunoo}{
  morekeywords={true,false,null},
  keywordstyle=\color{lstEmph}\bfseries,
  sensitive=true,
  morestring=[b]"
}
\lstdefinelanguage{mermaid}{
  morekeywords={flowchart,graph,sequenceDiagram,erDiagram,classDiagram,gantt,
            subgraph,end,participant,actor,alt,else,loop,par,opt,note,activate,
            deactivate,direction,TB,LR,RL,BT,Person,Person_Ext,System,System_Ext,
            SystemDb,Container,ContainerDb,ContainerQueue,Component,ComponentDb,
            Rel,Rel_D,Rel_U,Rel_L,Rel_R,System_Boundary,Container_Boundary,
            Enterprise_Boundary,Boundary,Deployment_Node,UpdateLayoutConfig,
            UpdateRelStyle,classDef,class,style,state,C4Context,C4Container,
            C4Component,C4Deployment,stateDiagram,stateDiagram-v2},
  sensitive=true,
  morecomment=[l]{\%\%},
  morestring=[b]"
}
\lstdefinelanguage{bashunoo}[]{bash}{
  morekeywords={pnpm,docker,gcloud,node,psql,tectonic,turbo,prisma,gsutil},
  sensitive=true
}

% Codigo em linha.
%
% Uso obrigatorio de \texttt e \textcolor (e nao de {\ttfamily\color{...}}):
% os comandos \text... emitem \leavevmode antes de trocar a fonte, o que faz o
% paragrafo comecar na fonte do texto corrente. Sem isso, uma celula de tabela
% que inicia com codigo teria a primeira linha de base calculada com as metricas
% da fonte monoespacada e ficaria visivelmente desalinhada das demais colunas.
\newcommand{\cd}[1]{\texttt{\textcolor{unooNavy}{#1}}}
\newcommand{\arq}[1]{\texttt{\textcolor{unooInkSoft}{#1}}}
\newcommand{\rota}[1]{\texttt{\textcolor{unooGreen}{#1}}}
\newcommand{\env}[1]{\texttt{\textcolor{lstEmph}{#1}}}
% Ponto de quebra de linha sem hifen, para caminho longo em celula estreita:
% identificador monoespacado nao hifeniza, e sem isto ele transborda a caixa.
\newcommand{\B}{\discretionary{}{}{}}

% --- Titulos ----------------------------------------------------------------
\usepackage{titlesec}
\usepackage{titletoc}

\titleformat{\chapter}[display]
  {\normalfont}
  {\raggedleft\fontsize{62}{62}\selectfont\color{unooBlueLight}\textbf{\thechapter}}
  {-10pt}
  {\raggedright\fontsize{23}{27}\selectfont\bfseries\color{unooNavy}}
  [\vspace{5pt}{\color{unooBlue}\rule{\linewidth}{2.2pt}}]
\titlespacing*{\chapter}{0pt}{-30pt}{26pt}

\titleformat{name=\chapter,numberless}[display]
  {\normalfont}
  {}
  {0pt}
  {\raggedright\fontsize{23}{27}\selectfont\bfseries\color{unooNavy}}
  [\vspace{5pt}{\color{unooBlue}\rule{\linewidth}{2.2pt}}]

\titleformat{\section}
  {\normalfont\Large\bfseries\color{unooNavy}}{\thesection}{0.7em}{}
\titleformat{\subsection}
  {\normalfont\large\bfseries\color{unooInk}}{\thesubsection}{0.7em}{}
\titleformat{\subsubsection}
  {\normalfont\normalsize\bfseries\color{unooInkSoft}}{\thesubsubsection}{0.7em}{}
\titleformat{\paragraph}[runin]
  {\normalfont\normalsize\bfseries\color{unooNavy}}{}{0pt}{}[\;---\;]

\titlespacing*{\section}{0pt}{20pt}{8pt}
\titlespacing*{\subsection}{0pt}{15pt}{6pt}
\titlespacing*{\subsubsection}{0pt}{12pt}{5pt}
\titlespacing*{\paragraph}{0pt}{9pt}{0.4em}

\setcounter{secnumdepth}{3}
\setcounter{tocdepth}{2}

% --- Sumario ----------------------------------------------------------------
\titlecontents{chapter}[0em]
  {\addvspace{9pt}\bfseries\color{unooNavy}}
  {\contentslabel[\thecontentslabel]{1.5em}}
  {}
  {\hfill\contentspage}
\titlecontents{section}[1.7em]
  {\addvspace{1.5pt}}
  {\contentslabel{2.2em}}
  {}
  {\titlerule*[0.7pc]{\color{unooAxis}.}\contentspage}
\titlecontents{subsection}[4.0em]
  {\addvspace{0.5pt}\small\color{unooInkSoft}}
  {\contentslabel{2.8em}}
  {}
  {\titlerule*[0.7pc]{\color{unooGrid}.}\contentspage}

% --- Cabecalho e rodape -----------------------------------------------------
\usepackage{fancyhdr}
\usepackage{lastpage}

\newcommand{\docid}{PU-DOC-001}
\newcommand{\docversao}{1.2}

\renewcommand{\headrule}{{\color{unooGrid}\hrule height \headrulewidth}}
\renewcommand{\footrule}{{\color{unooGrid}\hrule height \footrulewidth}}

\fancypagestyle{unoomain}{%
  \fancyhf{}
  \fancyhead[L]{\footnotesize\color{unooInkSoft}\nouppercase{\leftmark}}
  \fancyhead[R]{\raisebox{-0.15ex}{\includegraphics[height=1.8ex]{figuras/marca/simbolo_azul_profundo.png}}\hspace{4pt}\footnotesize\bfseries\color{unooNavy}Unooria}
  \fancyfoot[L]{\footnotesize\color{unooMuted}\docid{} \textperiodcentered{} v\docversao}
  \fancyfoot[R]{\footnotesize\color{unooInkSoft}\thepage\,/\,\pageref{LastPage}}
  \renewcommand{\headrulewidth}{0.5pt}
  \renewcommand{\footrulewidth}{0.5pt}
}
% Pre-textual: numeracao romana, sem total de paginas (o total do
% \pageref{LastPage} refere-se a numeracao arabe do corpo do documento).
\fancypagestyle{unoofront}{%
  \fancyhf{}
  \fancyhead[L]{\footnotesize\color{unooInkSoft}\nouppercase{\leftmark}}
  \fancyhead[R]{\raisebox{-0.15ex}{\includegraphics[height=1.8ex]{figuras/marca/simbolo_azul_profundo.png}}\hspace{4pt}\footnotesize\bfseries\color{unooNavy}Unooria}
  \fancyfoot[L]{\footnotesize\color{unooMuted}\docid{} \textperiodcentered{} v\docversao}
  \fancyfoot[R]{\footnotesize\color{unooInkSoft}\thepage}
  \renewcommand{\headrulewidth}{0.5pt}
  \renewcommand{\footrulewidth}{0.5pt}
}
\fancypagestyle{plain}{%
  \fancyhf{}
  \fancyfoot[R]{\footnotesize\color{unooInkSoft}\thepage}
  \renewcommand{\headrulewidth}{0pt}
  \renewcommand{\footrulewidth}{0pt}
}
\pagestyle{unoofront}
\renewcommand{\chaptermark}[1]{\markboth{\thechapter\quad #1}{}}

% --- Referencias cruzadas e PDF --------------------------------------------
\usepackage[
  colorlinks=true,
  linkcolor=unooNavy,
  citecolor=unooBlue,
  urlcolor=unooBlue,
  filecolor=unooBlue,
  bookmarksnumbered=true,
  bookmarksopen=true,
  bookmarksopenlevel=1,
  pdfdisplaydoctitle=true
]{hyperref}
\usepackage{bookmark}
\usepackage{url}
\urlstyle{sf}

% --- Bibliografia (ABNT NBR 6023) ------------------------------------------
%
% `abntex2cite` + BibTeX, e nao biblatex+biber: o pipeline de compilacao
% (Tectonic) chama o BibTeX classico automaticamente, sem ferramenta externa.
% A saida sai formatada em ABNT do mesmo jeito.
\usepackage[alf]{abntex2cite}

% --- Semantica de referencia do projeto ------------------------------------
\newcommand{\rf}[1]{\textbf{\textcolor{unooNavy}{#1}}}
\newcommand{\adrr}[1]{\textbf{\textcolor{unooLavenderSoft}{#1}}}
\newcommand{\ucr}[1]{\textbf{\textcolor{unooBlue}{#1}}}
\newcommand{\drr}[1]{\textbf{\textcolor{unooOrange}{#1}}}
\newcommand{\yes}{\textcolor{unooGreen}{\textbf{\checkmark}}}
\newcommand{\no}{\textcolor{unooOrange}{\textbf{--}}}

% Setas: a Roboto NAO tem U+2192 nem U+21D2 (medido — o XeTeX avisa
% "Missing character" e o glifo simplesmente nao e desenhado). Em prosa,
% sempre estes comandos; nunca a seta literal.
\newcommand{\seta}{\ensuremath{\rightarrow}}
\newcommand{\implica}{\ensuremath{\Rightarrow}}
\newcommand{\setabi}{\ensuremath{\leftrightarrow}}

% Espacamento de paragrafo
\setlength{\parindent}{0pt}
\setlength{\parskip}{7pt}

\widowpenalty=10000
\clubpenalty=10000
\flushbottom
